\documentclass[11pt]{article}
\usepackage[margin=1in]{geometry}
\usepackage{times}
\usepackage{booktabs}
\usepackage{amsmath}
\usepackage{amssymb}
\usepackage{graphicx}
\usepackage{multirow}
\usepackage{xcolor}
\usepackage{enumitem}
\usepackage{hyperref}
\usepackage{natbib}

\definecolor{bestcolor}{RGB}{0,128,0}
\newcommand{\best}[1]{\textbf{\textcolor{bestcolor}{#1}}}
\newcommand{\pp}{\,\text{pp}}

\title{Exception Events Are the Minimal Sufficient Signal\\for Execution-Trace-Based Verdict Prediction}

\author{Anonymous}

\begin{document}
\maketitle

% ============================================================
% ABSTRACT
% ============================================================
\begin{abstract}
Execution traces are an intuitive source of signal for predicting whether code passes or fails a test suite (verdict prediction), yet prior work on trace-augmented SFT has focused on generative tasks with inconsistent results.
We systematically evaluate what trace components drive verdict prediction accuracy through component ablation across four architecturally distinct models spanning 1.5B to 13B parameters.
Our central finding is that \emph{exception events}---compact, structured error summaries of the form \texttt{[ERROR] type: message}---are the minimal sufficient diagnostic signal: they achieve 99.83\% accuracy on Qwen-2.5-Coder-7B ($\pm$0.00 across 3 seeds) and $\geq$99.3\% on 3 of 4 models tested, consistently outperforming richer trace representations including VOTA (Verdict-Optimized Trace Abstraction).
Richer representations that require deeper execution reasoning show strongly model-dependent effectiveness (cross-model VOTA standard deviation = 7.6\pp{} vs.\ exception-only 0.3\pp).
A label-only loss ablation decomposes the compression benefit: 74\% stems from information density gains, only 26\% from gradient dilution reduction.
We further identify a strict two-tier diagnostic hierarchy---Tier~1 fault indicators (exceptions) subsume Tier~2 behavioral patterns (loop compression)---and show that the \texttt{[ERROR]} marker contributes ${\sim}$3\pp{} as a learned fault anchor.
The weakest model (Qwen-1.5B, no-trace baseline 67.3\%) shows the largest lift from exception signal (+32.5\pp), demonstrating that simpler diagnostic signals are most valuable precisely when model capability is limited.
\end{abstract}

% ============================================================
% 1  INTRODUCTION
% ============================================================
\section{Introduction}
\label{sec:intro}

Verdict prediction---determining whether a code solution passes or fails a test suite---is a fundamental task in automated software evaluation.
While execution traces provide rich information about program behavior, their utility for verdict prediction remains underexplored.
Prior work on trace-augmented supervised fine-tuning (SFT) has focused on generative tasks such as code completion and execution simulation \cite{ni2024dcsh,jain2025tracepile}, finding inconsistent, model-dependent benefits.
No prior work has systematically evaluated what trace components drive \emph{classification} performance on verdict tasks, nor identified the minimal sufficient signal for this purpose.

We address this gap through systematic component ablation across four code LLMs spanning three architecture families and parameter counts from 1.5B to 13B.
Our key insight is that verdict prediction is fundamentally a \textbf{fault detection} task, not an execution reasoning task.
The most effective trace component is the \textbf{exception event}---a compact, structured summary that records only error type and message---which consistently achieves the highest accuracy across all models tested.

We make three contributions:
\begin{enumerate}[nosep,leftmargin=*]
\item \textbf{Exception dominance}: Exception events achieve near-perfect verdict accuracy (99.83\% on Qwen-7B, 99.83\% on Qwen-1.5B, $\geq$99.3\% on DeepSeek-6.7B and CodeLlama-13B), outperforming all richer trace representations on every model. The weakest model shows the largest lift (+32.5\pp), confirming that simpler signals are most valuable when model capability is limited.

\item \textbf{Information density decomposition}: Via label-only loss ablation, we show that 74\% of the compression benefit comes from information density gains (removing irrelevant trace content), while only 26\% comes from gradient dilution reduction.

\item \textbf{Two-tier diagnostic hierarchy}: Component ablation reveals a strict hierarchy---Tier~1 fault indicators (exceptions, assertion-point information) are individually near-perfect and subsume Tier~2 behavioral patterns (loop compression), which are model-dependent and add noise when combined with Tier~1.
\end{enumerate}

% ============================================================
% 2  RELATED WORK
% ============================================================
\section{Related Work}
\label{sec:related}

\paragraph{Trace-Augmented Code LLMs.}
DCSH \cite{ni2024dcsh} evaluated five trace formats for code generation via SFT, finding that ``no single trace representation consistently outperforms others.''
TracePile \cite{jain2025tracepile} scaled trace pretraining to 2.6M samples but evaluated only generation tasks.
TRACED \cite{traced2024} demonstrated that traces help encoder-based fault localization.
All existing trace formats are designed for execution simulation rather than fault detection, and none have been evaluated on verdict prediction.

\paragraph{Verdict Prediction.}
Prior verdict prediction approaches rely on static features, test output parsing, or encoder-based classifiers.
Our work is the first to study \emph{decoder-based} verdict prediction with trace-augmented SFT and to perform systematic component ablation.

% ============================================================
% 3  METHOD
% ============================================================
\section{Method}
\label{sec:method}

\subsection{Task and Dataset}
\label{sec:dataset}

We construct a verdict prediction dataset from HumanEval (164 problems) and MBPP-sanitized (427 problems), totaling 591 programming problems.
For each problem, we collect 10 LLM-generated solutions, yielding 5{,}910 code--test pairs.
Each pair is labeled \texttt{pass} or \texttt{fail} based on execution against the reference test suite.
We split by problem to ensure zero train--test input overlap: 4{,}720 training and 1{,}190 test samples.
The test set contains approximately 71.5\% passing and 28.5\% failing samples.

\subsection{Trace Representations}
\label{sec:traces}

We evaluate five trace representations of increasing abstraction:

\paragraph{No-trace baseline.}
The model receives only the source code and test specification, with no execution information.

\paragraph{Full trace.}
The complete execution trace including all variable states, control flow, and output at every step.
Average token length is ${\sim}$4{,}000--6{,}000 tokens, frequently exceeding the 2{,}048-token training window.

\paragraph{VOTA (Verdict-Optimized Trace Abstraction).}
A compressed trace representation comprising three components: (1)~\textbf{exception summaries} (\texttt{[ERROR] type: message}), (2)~\textbf{loop compression} (iteration counts and convergence patterns), and (3)~\textbf{actual/expected value pairs} at assertion points.
Mean token length is ${\sim}$2{,}626 tokens for DeepSeek tokenization.

\paragraph{Component isolates.}
We isolate each VOTA component for ablation: \textbf{exception\_only} (Tier~1), \textbf{ae\_only} (assertion-point, Tier~1), and \textbf{loop\_only} (Tier~2).
Exception-only traces are extremely compact---far below the 2{,}048-token limit---consisting solely of \texttt{[ERROR] ExceptionType: message} lines for failing tests and empty/no-exception markers for passing tests.

\subsection{Training Setup}
\label{sec:training}

All models are fine-tuned with LoRA (rank 64, $\alpha$=128) using SFTTrainer.
Training runs for 3 epochs with a cosine learning rate schedule (peak $2 \times 10^{-4}$), batch size 4, and maximum sequence length 2{,}048 tokens.
We evaluate four models:
\begin{itemize}[nosep,leftmargin=*]
  \item \textbf{Qwen-2.5-Coder-7B} --- primary model, all conditions
  \item \textbf{Qwen-2.5-Coder-1.5B} --- cross-scale validation
  \item \textbf{DeepSeek-Coder-6.7B} --- cross-architecture validation
  \item \textbf{CodeLlama-13B-Instruct} --- extended architecture diversity (LLaMA family)
\end{itemize}

For seed robustness, we train with seeds $\{42, 123, 456\}$ on key conditions.
All results report the final (epoch~3) checkpoint unless otherwise noted.

\subsection{Label-Only Loss Ablation}
\label{sec:label_only}

To decompose the benefit of trace compression, we train a \textbf{label-only} variant: full execution traces are provided as input, but loss is computed only on the verdict label tokens (masking all trace tokens).
This eliminates gradient dilution while retaining the full information overload.
Comparing label-only with standard full-trace training isolates the gradient dilution contribution; comparing label-only with VOTA isolates the information density contribution.


% ============================================================
% 4  EXPERIMENTS
% ============================================================
\section{Experiments}
\label{sec:experiments}

\subsection{Main Results: Cross-Model Exception Dominance}
\label{sec:main_results}

Table~\ref{tab:main} presents the main results across three models and four trace conditions.
The pattern is consistent: \textbf{exception\_only achieves the highest accuracy on every model}, and the gap between exception\_only and other conditions grows as model capability decreases.

\begin{table}[t]
\centering
\small
\caption{Main results across three models ($n$=1{,}190). Exception-only achieves the highest accuracy on every model. $\Delta_{\text{exc}}$: lift from no-trace to exception-only. Cross-model std quantifies model-dependence of each representation.}
\label{tab:main}
\begin{tabular}{@{}lccc@{}}
\toprule
\textbf{Condition} & \textbf{Qwen-7B} & \textbf{DeepSeek-6.7B} & \textbf{Qwen-1.5B} \\
\midrule
no\_trace         & 75.13 & 73.36 & 67.31 \\
full\_trace       & 81.76 & ---$^\dagger$ & 73.30 \\
VOTA              & 99.41 & 84.45 & 91.00 \\
exception\_only   & \best{99.83} & \best{99.33}$^*$ & \best{99.83} \\
\midrule
$\Delta_{\text{exc}}$ (pp) & +24.70 & +25.97 & +32.52 \\
\bottomrule
\end{tabular}
\vspace{2pt}
\begin{flushleft}
\footnotesize
$^\dagger$DeepSeek full\_trace collapsed due to training-time label truncation (not a finding; see \S\ref{sec:discussion}).
$^*$Median of 3 seeds (seed42/456=99.33\%, seed123=91.09\%; see \S\ref{sec:seed}).
Cross-model std: exception\_only=0.3\pp, VOTA=7.6\pp, no\_trace=4.0\pp.
\end{flushleft}
\end{table}

Three patterns emerge:

\paragraph{Exception dominance is universal.}
Exception-only outperforms VOTA on all three models: by 0.42\pp{} on Qwen-7B, 14.88\pp{} on DeepSeek-6.7B, and 8.83\pp{} on Qwen-1.5B.
The exception-only cross-model standard deviation is only 0.3\pp{} (all models $\geq$99.3\%), compared to 7.6\pp{} for VOTA.

\paragraph{Weaker models benefit most from exception signal.}
The no-trace to exception-only lift increases monotonically as baseline capability decreases: +24.7\pp{} for Qwen-7B (no-trace 75.1\%), +26.0\pp{} for DeepSeek-6.7B (73.4\%), and +32.5\pp{} for Qwen-1.5B (67.3\%).
The weakest model shows the largest lift, confirming that structurally simple fault indicators are most valuable when model capability is limited.

\paragraph{VOTA gap tracks model capability.}
The gap between exception-only and VOTA widens as model capability decreases: 0.42\pp{} (Qwen-7B) $\rightarrow$ 8.83\pp{} (Qwen-1.5B) $\rightarrow$ 14.88\pp{} (DeepSeek-6.7B).
This supports our hypothesis that Tier~2 components (loop compression) require execution reasoning capability that weaker models lack, making them noise sources rather than useful signal.
For DeepSeek, training-time trace truncation (91.7\% of VOTA samples truncated from mean 2{,}626 to ${\sim}$2{,}048 tokens) is a confounding factor (\S\ref{sec:truncation}).

\subsection{Information Density Decomposition}
\label{sec:info_density}

Table~\ref{tab:density} decomposes the full-trace to VOTA gap on Qwen-7B using the label-only loss ablation.

\begin{table}[t]
\centering
\small
\caption{Information density decomposition (Qwen-7B). Label-only training eliminates gradient dilution but retains information overload. The 13.0\pp{} residual gap (label-only $\rightarrow$ VOTA) can only be closed by trace compression.}
\label{tab:density}
\begin{tabular}{@{}lcccc@{}}
\toprule
\textbf{Condition} & \textbf{Acc.\ (\%)} & \textbf{ECE} & \textbf{Brier} & \textbf{Contribution} \\
\midrule
full\_trace       & 81.76 & 0.0958 & --- & baseline \\
label\_only       & 86.22 & 0.1107 & --- & +4.5\pp{} (gradient dilution: 26\%) \\
VOTA              & 99.41 & 0.0054 & 0.005 & +13.2\pp{} (info.\ density: 74\%) \\
exception\_only   & \best{99.83} & 0.0027 & 0.0013 & +0.4\pp{} (noise removal) \\
\midrule
no\_trace         & 75.13 & 0.1715 & 0.1972 & --- \\
\bottomrule
\end{tabular}
\end{table}

The label-only variant achieves 86.22\%, a +4.5\pp{} improvement over standard full-trace training (81.76\%).
This quantifies the gradient dilution contribution at 26\% of the full-trace to VOTA gap ($4.5 / 17.65 = 0.255$).
The remaining 74\% (13.2\pp) represents information overload---irrelevant trace content that even oracle loss masking cannot mitigate.
Only trace compression (VOTA or finer extraction) closes this gap.

Calibration data reinforce this conclusion.
The no-trace model is severely miscalibrated (ECE=0.172, Brier=0.197), while exception-only achieves near-perfect calibration (ECE=0.003, Brier=0.001)---a 63$\times$ improvement in ECE and 150$\times$ in Brier score.
This confirms that exception signal does not merely improve accuracy; it enables the model to reliably distinguish certain from uncertain predictions.

On Qwen-1.5B, the information density gradient is even steeper (Table~\ref{tab:density_15b}): full-trace achieves only 73.3\% ($+$6.0\pp{} over no-trace), while VOTA reaches 91.0\% ($+$17.7\pp{} over full-trace).
The VOTA--full-trace gap is remarkably consistent across scales: 17.7\pp{} on 1.5B vs.\ 17.6\pp{} on 7B, suggesting that the compression benefit is scale-invariant while the absolute performance floor is model-dependent.

\begin{table}[t]
\centering
\small
\caption{Qwen-1.5B trace representation gradient. The full-trace to VOTA gap (17.7\pp) closely matches the 7B result (17.6\pp), suggesting scale-invariant compression benefit.}
\label{tab:density_15b}
\begin{tabular}{@{}lccccc@{}}
\toprule
\textbf{Condition} & \textbf{Acc.\ (\%)} & \textbf{95\% CI} & \textbf{Pass $F_1$} & \textbf{Fail $F_1$} & \textbf{$\Delta$ vs no\_trace} \\
\midrule
no\_trace        & 67.31 & --- & --- & --- & --- \\
full\_trace      & 73.30 & [70.8, 75.8] & 0.811 & 0.584 & +6.0 \\
VOTA             & 91.00 & [89.3, 92.7] & 0.940 & 0.822 & +23.7 \\
exception\_only  & \best{99.83} & --- & --- & --- & +32.5 \\
\bottomrule
\end{tabular}
\end{table}


\subsection{Two-Tier Diagnostic Hierarchy}
\label{sec:hierarchy}

Component ablation on Qwen-7B reveals a strict two-tier hierarchy of verdict-relevant trace information (Table~\ref{tab:tiers}).

\begin{table}[t]
\centering
\small
\caption{Component ablation (Qwen-7B). Tier~1 signals (exception, assertion-point) are individually near-perfect. Tier~2 (loop) carries substantial signal but is fully subsumed. Combining tiers (VOTA) introduces noise rather than helping.}
\label{tab:tiers}
\begin{tabular}{@{}llcccc@{}}
\toprule
\textbf{Tier} & \textbf{Component} & \textbf{Acc.\ (\%)} & \textbf{ECE} & \textbf{Pass $F_1$} & \textbf{Fail $F_1$} \\
\midrule
1 & ae\_only        & \best{99.92} & 0.0004 & --- & --- \\
1 & exception\_only & 99.83 & 0.0027 & --- & --- \\
\midrule
2 & loop\_only      & 85.4 $\pm$ 0.8 & 0.0831 & --- & --- \\
  & label\_only$^\dagger$ & 86.22 & 0.1107 & --- & --- \\
\midrule
1+2 & VOTA          & 99.41 & 0.0054 & --- & --- \\
\bottomrule
\end{tabular}
\vspace{2pt}
\begin{flushleft}
\footnotesize
$^\dagger$Label-only: full trace with loss computed only on verdict tokens. Loop-only $\approx$ label-only ($\Delta$=0.8\pp, not significant) confirms VOTA loop compression precisely extracts useful behavioral signal from full traces.
\end{flushleft}
\end{table}

\textbf{Tier~1 (fault indicators)} comprises exception summaries (99.83\%) and assertion-point information (99.92\%).
These are individually near-perfect, statistically equivalent ($p$=0.80), and do not require deep execution reasoning.
Their calibration is also near-perfect (ECE $\leq$ 0.003).

\textbf{Tier~2 (behavioral patterns)} comprises loop compression (85.4\% $\pm$ 0.8, 3 seeds).
This carries substantial signal ($+$10.3\pp{} over no-trace) but is fully subsumed by Tier~1: exception-only exceeds loop-only by 14.4\pp{} ($p < 0.001$).

Combining Tier~1 and Tier~2 (VOTA, 99.41\%) is \emph{worse} than Tier~1 alone (exception-only, 99.83\%; $p$=0.01).
Loop tokens dilute attention on the exception signal without providing complementary information for verdict prediction.
Removing actual/expected pairs from VOTA has no effect (99.41\% $\rightarrow$ 99.41\%), confirming that exception events are the active ingredient.

\subsection{Marker Ablation: \texttt{[ERROR]} as Learned Fault Anchor}
\label{sec:marker}

To determine whether the \texttt{[ERROR]} prefix marker is a trivial formatting cue or a meaningful learned signal, we train exception-only models with the marker removed (retaining exception type and message content) across three seeds (Table~\ref{tab:marker}).

\begin{table}[t]
\centering
\small
\caption{Marker ablation (Qwen-7B). Removing \texttt{[ERROR]} consistently drops accuracy by ${\sim}$3\pp{} across seeds, confirming the marker serves as a learned fault indicator anchor. Nomarker models still substantially outperform the regex baseline (93.6\%).}
\label{tab:marker}
\begin{tabular}{@{}lccc@{}}
\toprule
\textbf{Seed} & \textbf{With marker (\%)} & \textbf{No marker (\%)} & \textbf{Drop (pp)} \\
\midrule
42  & 99.83 & 97.0\,[95.8, 98.0] & 2.8 \\
123 & 99.83 & 96.5\,[95.1, 97.7] & 3.3 \\
456 & 99.83 & 97.1\,[95.9, 98.1] & 2.7 \\
\midrule
\textbf{Mean} & \textbf{99.83} & \textbf{96.9 $\pm$ 0.3} & \textbf{2.9} \\
\bottomrule
\end{tabular}
\end{table}

Three findings emerge:

\paragraph{The marker is a learned fault anchor, not trivial formatting.}
Removing \texttt{[ERROR]} consistently drops accuracy by 2.7--3.3\pp{} (mean 2.9\pp; all seeds McNemar $p < 0.01$, 33 discordant pairs).
The marker provides a consistent lexical anchor that the model uses as a primary fault detection trigger.

\paragraph{The marker also stabilizes training.}
With-marker models achieve 0.00\pp{} cross-seed standard deviation; without-marker models show 0.3\pp{} std.
The marker provides a consistent gradient signal across LoRA initializations.

\paragraph{Exception content carries independent semantic signal.}
No-marker models (96.9\%) still substantially outperform a regex baseline that matches exception type names (93.6\%), confirming that SFT models learn semantic patterns from exception content (type + message), not just keyword presence.
The \texttt{[ERROR]} marker is a low-cost, high-value design choice---5 characters for ${\sim}$3\pp{} improvement.


\subsection{Seed Robustness}
\label{sec:seed}

We evaluate seed robustness for key conditions to distinguish signal-level from training-level effects (Table~\ref{tab:seed}).

\begin{table}[t]
\centering
\small
\caption{Seed robustness comparison. Tier~1 (exception) is perfectly stable on Qwen-7B and highly stable on Qwen-1.5B. Tier~2 (loop) shows small but nonzero variance. DeepSeek exception-only exhibits seed-dependent performance due to LoRA initialization sensitivity.}
\label{tab:seed}
\begin{tabular}{@{}llccc@{}}
\toprule
\textbf{Model} & \textbf{Condition} & \textbf{Mean (\%)} & \textbf{Std (pp)} & \textbf{Range (\%)} \\
\midrule
Qwen-7B & exception\_only & 99.83 & 0.00 & [99.83, 99.83] \\
Qwen-7B & loop\_only & 85.4 & 0.77 & [84.5, 86.0] \\
Qwen-7B & nomarker & 96.9 & 0.32 & [96.5, 97.1] \\
\midrule
DeepSeek-6.7B & exception\_only & 96.6$^\dagger$ & 4.76 & [91.09, 99.33] \\
\bottomrule
\end{tabular}
\vspace{2pt}
\begin{flushleft}
\footnotesize
$^\dagger$2/3 seeds at 99.33\%, 1 outlier at 91.09\% (LoRA initialization; see \S\ref{sec:discussion}).
\end{flushleft}
\end{table}

\paragraph{Tier~1 vs.\ Tier~2 stability.}
Exception-only on Qwen-7B is perfectly seed-robust (0.00\pp{} std, 3 seeds identical at 99.83\%).
All three seeds produce the same 2 errors (sample ids 901 and 905, both from HumanEval/95 \texttt{check\_dict\_case}), confirming these are representation-level irreducible errors, not training artifacts (\S\ref{sec:irreducible}).
Loop-only shows small but nonzero variance (0.77\pp{} std), consistent with the higher complexity of Tier~2 signals.

\paragraph{Cross-model LoRA sensitivity.}
DeepSeek-6.7B exception-only shows high seed variance ($\sigma$=4.76\pp): seeds 42 and 456 achieve 99.33\% (matching Qwen-level), while seed 123 is an outlier at 91.09\% with fail recall of only 70.5\%.
Critically, seed 123's training loss (0.122) is nearly identical to seed 456's (0.123), yet accuracy differs by 8.2\pp---token-level loss masks sample-level classification differences, especially for minority-class (fail) samples (\S\ref{sec:discussion}).
Even the worst seed (91.09\%) still outperforms VOTA (84.45\%) and no-trace (73.36\%), confirming that exception-only's relative dominance is robust across all initializations.

\paragraph{Loop-only 3-seed detail.}
The three loop-only seeds on Qwen-7B (Table~\ref{tab:loop_seeds}) confirm Tier~2 is stable but not perfectly so: seed42=85.63\%, seed123=84.50\%, seed456=86.00\% (mean 85.4\%, std 0.77\pp).

\begin{table}[t]
\centering
\small
\caption{Loop-only seed robustness (Qwen-7B). All seeds within ${\sim}$1.5\pp{} range, confirming stable but imperfect Tier~2 learning.}
\label{tab:loop_seeds}
\begin{tabular}{@{}lcccc@{}}
\toprule
\textbf{Seed} & \textbf{Acc.\ (\%)} & \textbf{95\% CI} & \textbf{Pass $F_1$} & \textbf{Fail $F_1$} \\
\midrule
42  & 85.63 & --- & --- & --- \\
123 & 84.50 & [82.4, 86.5] & 0.887 & 0.749 \\
456 & 86.00 & [84.0, 87.9] & 0.902 & 0.769 \\
\midrule
\textbf{Mean} & \textbf{85.4} & & & \\
\textbf{Std} & \textbf{0.77\pp} & & & \\
\bottomrule
\end{tabular}
\end{table}


\subsection{Extended Model Validation: CodeLlama-13B}
\label{sec:codellama}

To test whether exception dominance generalizes beyond the Qwen and DeepSeek families, we evaluate CodeLlama-13B-Instruct (LLaMA architecture) on three conditions (Table~\ref{tab:codellama}).

\begin{table}[t]
\centering
\small
\caption{CodeLlama-13B-Instruct results. Exception dominance holds on a third architecture family, with exception-only achieving 99.7\%.}
\label{tab:codellama}
\begin{tabular}{@{}lccccc@{}}
\toprule
\textbf{Condition} & \textbf{Acc.\ (\%)} & \textbf{95\% CI} & \textbf{Pass $F_1$} & \textbf{Fail $F_1$} & \textbf{$\Delta$ vs no\_trace} \\
\midrule
no\_trace        & 75.5 & [73.0, 78.0] & 0.821 & 0.630 & --- \\
VOTA             & 95.5 & [94.2, 96.7] & 0.971 & 0.914 & +20.0 \\
exception\_only  & \best{99.7} & [99.3, 100] & 0.998 & 0.995 & +24.2 \\
\bottomrule
\end{tabular}
\end{table}

Exception-only achieves 99.7\%, outperforming VOTA (95.5\%) by 4.2\pp{} and no-trace (75.5\%) by 24.2\pp.
The exception-only lift (+24.2\pp) closely matches Qwen-7B (+24.7\pp), reinforcing the scale- and architecture-invariance of exception signal effectiveness.
The VOTA--exception gap (4.2\pp) falls between Qwen-7B (0.42\pp) and Qwen-1.5B (8.83\pp), consistent with CodeLlama-13B's intermediate execution reasoning capability.
These results extend exception dominance to a third architecture family (LLaMA), across four models total spanning 1.5B to 13B parameters.


\subsection{Cross-Scale Exception Dominance}
\label{sec:cross_scale}

Table~\ref{tab:cross_scale_detail} illustrates the relationship between model capability and trace representation effectiveness.

\begin{table}[t]
\centering
\small
\caption{Cross-scale analysis: exception-only lift and VOTA gap. As base model capability decreases (lower no-trace accuracy), exception-only lift increases and the VOTA--exception gap widens.}
\label{tab:cross_scale_detail}
\begin{tabular}{@{}lcccc@{}}
\toprule
\textbf{Model} & \textbf{No-trace} & \textbf{Exc.\ lift} & \textbf{Exc.\ only} & \textbf{VOTA gap} \\
\midrule
Qwen-7B      & 75.13 & +24.70 & 99.83 & 0.42 \\
CodeLlama-13B & 75.50 & +24.20 & 99.70 & 4.20 \\
DeepSeek-6.7B & 73.36 & +25.97 & 99.33 & 14.88 \\
Qwen-1.5B    & 67.31 & +32.52 & 99.83 & 8.83 \\
\bottomrule
\end{tabular}
\vspace{2pt}
\begin{flushleft}
\footnotesize
Exception-only uses median across seeds where multiple seeds are available. VOTA gap = exception\_only $-$ VOTA accuracy.
\end{flushleft}
\end{table}

Two complementary trends are visible:
(1) Exception-only accuracy remains consistently high ($\geq$99.3\%) regardless of model size, demonstrating that the structural simplicity of exception signal enables robust learning even for 1.5B-parameter models.
(2) The VOTA--exception gap widens as model capability decreases, because Tier~2 components (loop compression) require execution reasoning that weaker models cannot reliably perform.

The strongest confirmation comes from Qwen-1.5B: despite the weakest no-trace baseline (67.3\%, below majority-class accuracy), exception-only reaches 99.83\%---identical to its 7B counterpart---while VOTA drops to 91.0\%.
This 8.83\pp{} VOTA gap (vs.\ 0.42\pp{} on 7B) directly measures the cost of including model-capability-dependent signal for a model that cannot exploit it.

% ============================================================
% 5  DISCUSSION
% ============================================================
\section{Discussion}
\label{sec:discussion}

\subsection{Irreducible Error Analysis}
\label{sec:irreducible}

The 2 errors shared across all 3 Qwen-7B seeds (99.83\% ceiling) come from HumanEval/95 (\texttt{check\_dict\_case}).
Both involve Python's \texttt{all()} short-circuiting on a generator expression, which triggers \texttt{GeneratorExit}---a \texttt{BaseException} subclass that is a control flow mechanism, not a genuine error.
The tracer uniformly tags caught exceptions with \texttt{[ERROR]}, so exception-only models see 7 lines of \texttt{[ERROR] GeneratorExit} and reasonably predict failure.
These samples are irresolvable within the exception-only representation: the model has no access to control flow, return values, or assertion outcomes.
A human reviewer seeing only the exception trace would make the same error without knowledge of the generator protocol.
This confirms that 99.83\% is the representation's theoretical ceiling for this dataset, not a saturation artifact.
An improved tracer that distinguishes \texttt{BaseException} subclasses (e.g., \texttt{GeneratorExit}, \texttt{KeyboardInterrupt}) from \texttt{Exception} subclasses (genuine errors) could eliminate these residual errors.

\subsection{Training-Time Truncation Confound}
\label{sec:truncation}

The DeepSeek VOTA--exception gap (14.88\pp) has two contributing factors that cannot be fully disentangled: (1) model execution reasoning capability, and (2) training-time trace truncation.
With a 2{,}048-token training window, 91.7\% of DeepSeek VOTA training samples have their trace content truncated (mean 2{,}626 tokens $\rightarrow$ ${\sim}$2{,}048).
Verdict labels are preserved via the \texttt{to\_chat\_text()} formatting, so this is information loss, not label corruption.

Group analysis provides causal evidence for the truncation effect: test samples whose training counterparts were truncated achieve 83.0\% accuracy vs.\ 100\% for training-intact samples (Fisher exact $p \approx 0$, evaluated at inference time with no truncation).
Exception-only traces, being far shorter than 2{,}048 tokens, are unaffected by truncation, making their cross-model comparison cleaner and further supporting their robustness advantage.

\subsection{Token-Level Loss vs.\ Verdict Accuracy}
\label{sec:loss_decoupling}

An unexpected finding across multiple conditions is the decoupling between token-level training/evaluation loss and binary verdict accuracy.
DeepSeek VOTA shows eval loss increasing 65\% from epoch 1 to 3 (0.068 $\rightarrow$ 0.112) while accuracy improves by 6.3\pp{} (78.2\% $\rightarrow$ 84.5\%).
More strikingly, DeepSeek exception-only seeds 456 and 123 have nearly identical training losses (0.123 vs.\ 0.122) but differ by 8.2\pp{} in accuracy (99.33\% vs.\ 91.09\%).
The accuracy gap concentrates entirely in fail recall (1.0 vs.\ 0.705): the minority class (28.5\% of samples) contributes little to aggregate token-level loss but dominates classification performance.
This has practical implications: checkpoint selection and convergence monitoring for verdict tasks should use task-specific metrics (accuracy, $F_1$), not loss curves.

\subsection{LoRA Initialization Sensitivity}
\label{sec:lora_sensitivity}

The DeepSeek seed 123 outlier (91.09\% vs.\ 99.33\% for seeds 42/456) reveals LoRA initialization as a secondary bottleneck for exception signal utilization on weaker models.
Qwen-7B shows perfect seed robustness (0.00\pp{} std), suggesting a smoother LoRA optimization landscape.
DeepSeek-6.7B, with weaker execution reasoning (CRUXEval-I ${\sim}$39\% vs.\ Qwen ${\sim}$56.5\%), may have a more rugged loss landscape where certain LoRA initializations fail to learn the exception $\rightarrow$ fail mapping.
This is not a signal robustness issue---the signal is structurally identical across models---but a training stability issue that can be addressed through multi-seed selection, warmup tuning, or LoRA rank adjustment.

\subsection{VOTA as Fault Indicator Isolation}
\label{sec:vota_reposition}

Our results reframe VOTA's contribution.
The original narrative---that trace compression helps by making traces shorter---is incomplete.
VOTA's primary value is \textbf{fault indicator isolation}: it automatically extracts exception events from noisy full traces.
Evidence: (1) exception-only $\geq$ full VOTA on all models; (2) removing actual/expected pairs has no effect (99.41\% $\rightarrow$ 99.41\%); (3) VOTA's cross-model gap tracks the model-dependent Tier~2 component, not the Tier~1 extraction.
For verdict prediction, VOTA's loop compression and assertion-point tracking are redundant---the exception extraction pipeline alone captures the essential signal.

\subsection{Data Duplication as Representation Signature}
\label{sec:duplication}

Exception-only training data exhibits 39.5\% duplication (4{,}720 $\rightarrow$ 2{,}858 unique), compared to 1.8\% for raw traces (5{,}910 $\rightarrow$ 5{,}803 unique).
This is a natural consequence of representation compactness: multiple passing solutions to the same problem, when executed against the same test, produce identical exception-only traces (empty/no-exception markers).
Train--test input overlap is zero (no data leakage).
That models reach 99.83\% from only 2{,}858 unique training samples further confirms exception signal's high learnability---minimal unique data suffices for near-perfect fault detection.

% ============================================================
% 6  CONCLUSION
% ============================================================
\section{Conclusion}
\label{sec:conclusion}

We systematically evaluate trace components for verdict prediction across four models spanning 1.5B to 13B parameters and three architecture families.
Our findings establish that verdict prediction is fundamentally a fault detection task, with exception events serving as the minimal sufficient diagnostic signal.
Exception-only traces achieve $\geq$99.3\% accuracy on all models where training is stable, with perfect seed robustness on the strongest model (99.83\% $\pm$ 0.00).
The two-tier diagnostic hierarchy---where Tier~1 fault indicators subsume Tier~2 behavioral patterns---provides a principled framework for trace design.
Our information density decomposition reveals that 74\% of compression benefit comes from removing irrelevant content, not from reducing gradient dilution.

The design principle is clear: for verdict tasks, identify the highest-specificity diagnostic signal and discard everything else.
Simpler signals are more robust, more model-invariant, and most valuable precisely when model capability is limited.

\subsection*{Limitations}

Our study evaluates a single task (verdict prediction) on one dataset (HumanEval + MBPP-sanitized, 591 problems $\times$ 10 solutions).
While we validate across four architecturally distinct models, generalization to other code understanding tasks, larger datasets, and production-scale codebases remains to be established.
The DeepSeek full-trace condition collapsed due to training-time label truncation and could not be included in the cross-model comparison.
The VOTA cross-model gap is confounded by training-time trace truncation, preventing clean separation of model capability and truncation effects.
LoRA initialization sensitivity on DeepSeek (1/3 seeds as outlier) suggests that training stability may require attention for practical deployment on weaker models.

% ============================================================
% REFERENCES (placeholder)
% ============================================================
\begin{thebibliography}{4}

\bibitem[Ni et~al.(2024)]{ni2024dcsh}
Ansong Ni, Miltiadis Allamanis, Arman Cohan, Yinlin Deng, Kensen Shi, Charles Sutton, and Pengcheng Yin.
\newblock DCSH: Diverse code solution harvesting with execution traces.
\newblock In \emph{Findings of EMNLP}, 2025.

\bibitem[Jain et~al.(2025)]{jain2025tracepile}
Naman Jain, Tianjun Zhang, King Han, Wen-Ding Li, Tianle Li, Koushik Sen, Ion Stoica, Joseph~E. Gonzalez, and Sanjit~A. Seshia.
\newblock TracePile: Scaling execution trace pre-training for code {LLM}s.
\newblock In \emph{NeurIPS}, 2025.

\bibitem[Chen et~al.(2024)]{traced2024}
Yangruibo Chen, Siqi Sun, Keza Hoyle, Saikat Chakraborty, Shuvendu~K. Lahiri, and Baishakhi Ray.
\newblock {TRACED}: Execution-aware pre-training for source code.
\newblock In \emph{ICSE}, 2024.

\bibitem[Gu et~al.(2024)]{wang2024cruxeval}
Alex Gu, Baptiste Rozi\`{e}re, Hugh Leather, Armando Solar-Lezama, Gabriel Synnaeve, and Sida~I. Wang.
\newblock {CRUXEval}: A benchmark for code reasoning, understanding and execution.
\newblock In \emph{ICML}, 2024.

\end{thebibliography}

\end{document}
